\documentclass[11pt,a4paper]{article}

\usepackage[utf8]{inputenc}
\usepackage[T5]{fontenc} % For Vietnamese characters
\usepackage{amsmath, amssymb}
\usepackage{graphicx}
\usepackage{hyperref}
\usepackage{booktabs}
\usepackage{geometry}
\geometry{a4paper, margin=2.5cm}
\usepackage{enumitem}
\usepackage{float}
\usepackage{caption}
\usepackage{multirow}
\usepackage{tikz}
\usetikzlibrary{shapes.geometric, arrows.meta, positioning, fit, backgrounds}

\title{\textbf{Kiến trúc Causal GNN v2 cho Phân loại Thảm họa Đa phương thức: Phân tách Đặc trưng và Điều chỉnh Cửa sau}}
\author{Nhóm Nghiên cứu Dự án CrisisSummarization}
\date{\today}

\begin{document}

\maketitle

\begin{abstract}
Việc phân loại dữ liệu thảm họa từ mạng xã hội đóng vai trò quan trọng trong việc hỗ trợ các nỗ lực cứu trợ nhân đạo. Tuy nhiên, các mô hình học sâu truyền thống thường bị ảnh hưởng bởi những sự tương quan nhiễu (spurious correlations) do thiên lệch bối cảnh (domain bias). Báo cáo này trình bày kiến trúc Causal GNN v2, một phương pháp học sâu nhằm cải thiện khả năng tổng quát hóa trên miền dữ liệu thảm họa bằng cách tích hợp suy luận nhân quả (Causal Inference) và Mạng nơ-ron Tích chập Đồ thị (Graph Convolutional Networks - GCN). Thông qua việc phân tách đặc trưng (Disentanglement) thành nhánh nhân quả và nhánh nhiễu, đồng thời sử dụng thuật toán Điều chỉnh Cửa sau (Backdoor Adjustment) với cơ chế Monte Carlo, mô hình loại bỏ các yếu tố nhiễu trong quá trình suy luận. Thử nghiệm trên 3 tác vụ phân loại của bộ dữ liệu CrisisMMD cho thấy mô hình đạt độ chính xác cân bằng (Balanced Accuracy) và Weighted F1 cao, thiết lập hiệu suất cạnh tranh cho bài toán phân loại thảm họa.
\end{abstract}

\section{Giới thiệu (Introduction)}

Dữ liệu thảm họa trên mạng xã hội thường mang tính đa phương thức, kết hợp văn bản và hình ảnh. Mặc dù các mô hình nền tảng như CLIP cho khả năng trích xuất đặc trưng tốt, chúng dễ dàng "học vẹt" bối cảnh thay vì học nguyên nhân thực sự (causal factors) của thảm họa. Ví dụ, sự hiện diện của áo phao có thể gắn với lũ lụt, nhưng cũng có thể xuất hiện trong bối cảnh du lịch bình thường.

Các phương pháp phân loại SOTA hiện tại (như CAMO, SimMMDG, CIRL) mặc dù đã nỗ lực giải quyết phân phối ngoài miền (Out-of-Distribution) nhưng vẫn gặp hạn chế. Phần lớn chỉ dừng lại ở việc tách biệt đặc trưng cơ bản (Modality-level) hoặc chỉ giải quyết nhân quả trên dữ liệu đơn phương thức (Single-modal Graph).

Để giải quyết triệt để rào cản này, kiến trúc Causal GNN v2 đề xuất một cơ chế can thiệp nhân quả (Causal Intervention) thông qua lý thuyết đồ thị và tách bạch 3 cấp độ (3-level disentanglement). Thay vì dự đoán dựa vào mô hình nhân quả quan sát $\mathbb{P}(Y|X)$, chúng ta định hình lại dữ liệu qua toán tử \textit{do-calculus} của Judea Pearl: $\mathbb{P}(Y|do(X))$. Báo cáo này đi sâu vào phân tích kiến trúc, pipeline xử lý, định thức toán học và đánh giá toàn diện.
\section{Kiến trúc Hệ thống (System Architecture)}

Kiến trúc giải quyết bài toán đa phương thức thông qua 4 module trọng tâm, tạo thành một \textbf{Pipeline khép kín} từ dữ liệu thô đến dự đoán miễn nhiễm nhiễu.

\begin{figure}[H]
    \centering
    \resizebox{\linewidth}{!}{
    \begin{tikzpicture}[
        node distance=1.2cm and 1.5cm,
        box/.style={rectangle, draw=blue!70, fill=blue!5, thick, rounded corners, align=center, minimum height=0.8cm, font=\small},
        causal/.style={rectangle, draw=green!70, fill=green!5, thick, rounded corners, align=center, minimum height=0.8cm, font=\small},
        spurious/.style={rectangle, draw=red!70, fill=red!5, thick, rounded corners, align=center, minimum height=0.8cm, font=\small},
        graph/.style={rectangle, draw=purple!70, fill=purple!5, thick, rounded corners, align=center, minimum height=0.8cm, font=\small},
        arrow/.style={-{Stealth[scale=1.0]}, thick}
    ]

    % Nodes - Input & Projection
    \node[box] (img) {Image ($v$)};
    \node[box, below=1cm of img] (txt) {Text ($t$)};
    
    \node[box, right=of img] (proj_img) {ModalityProj \\ $z_{img\_g}, z_{img\_s}$};
    \node[box, right=of txt] (proj_txt) {ModalityProj \\ $z_{txt\_g}, z_{txt\_s}$};
    
    \node[box, right=of proj_img, yshift=-0.9cm] (unified) {Concat / Fusion \\ ($z_{unified} + \text{MixUp}$)};
    
    % Disentanglement
    \node[causal, right=of unified, yshift=1.4cm] (xc) {Causal Branch \\ $X_c$};
    \node[spurious, right=of unified, yshift=-1.4cm] (xs) {Spurious Branch \\ $X_s$};
    
    % Graph & Memory
    \node[box, above=0.5cm of xc, fill=yellow!10, draw=yellow!70] (hetero) {Heterophily Scorer \\ (Clean edges $\tilde{A}$)};
    \node[graph, right=of xc] (sage) {GraphSAGE Layer \\ $\text{ReLU}(W_{sage}[X_c || \tilde{A}X_c])$};
    \node[spurious, right=of xs] (bank) {Memory Bank \\ (Size: 2000)};
    
    % Causal Inference (Backdoor)
    \node[box, right=of sage, yshift=-1.4cm, fill=orange!10, draw=orange!70] (combine) {$\oplus$ Fusion \\ $X_c^{graph} + x_s^{(i)}$};
    \node[box, right=of combine, fill=orange!10, draw=orange!70] (classifier) {Classifier \& \\ Softmax};
    \node[box, right=of classifier, fill=orange!10, draw=orange!70] (output) {Robust Logits \\ $\log(\frac{1}{N}\sum P_i)$};

    % Arrows
    \draw[arrow] (img) -- (proj_img);
    \draw[arrow] (txt) -- (proj_txt);
    
    \draw[arrow] (proj_img.east) -- (unified.west);
    \draw[arrow] (proj_txt.east) -- (unified.west);
    
    \draw[arrow] (unified.east) -- (xc.west);
    \draw[arrow] (unified.east) -- (xs.west);
    
    \draw[arrow] (xc) -- (sage);
    \draw[arrow] (xc) -- (hetero);
    \draw[arrow] (hetero.east) -| (sage.top) node[pos=0.25, above, font=\scriptsize] {$\tilde{A}$};
    
    \draw[arrow] (xs) -- (bank) node[midway, above, font=\scriptsize, text=red] {Store};
    
    \draw[arrow] (sage.south) |- (combine.west) node[near start, right, font=\scriptsize] {$X_c^{graph}$};
    \draw[arrow, dashed] (bank.east) -| (combine.south) node[near start, below, font=\scriptsize, text=red] {Sample N=50 ($x_s^{(i)}$)};
    
    \draw[arrow] (combine) -- (classifier) node[midway, above, font=\scriptsize] {$\forall i \in [1, N]$};
    \draw[arrow] (classifier) -- (output) node[midway, above, font=\scriptsize] {$\mathbb{E}$};
    
    % Grouping
    \begin{scope}[on background layer]
        \node[fit=(proj_img)(proj_txt)(unified)(xc)(xs), draw=gray, dashed, inner sep=0.3cm, rounded corners, label=above:3-Level Disentanglement] {};
        \node[fit=(combine)(classifier)(output), draw=orange, dashed, inner sep=0.3cm, rounded corners, label=above:Causal Inference (Backdoor Adjustment)] {};
    \end{scope}

    \end{tikzpicture}
    }
    \caption{Sơ đồ luồng dữ liệu (Pipeline) của Kiến trúc Causal GNN v2}
    \label{fig:architecture}
\end{figure}

\subsection{Trích xuất và Phân tách Đặc trưng (Disentanglement)}
Đầu vào hệ thống gồm hình ảnh $\mathcal{I}$ và văn bản $\mathcal{T}$. Mô hình sử dụng bộ mã hóa \texttt{open\_clip} (ViT, Text Encoder) đã được pre-train:
\begin{align}
v &= \text{CLIP}_{vision}(\mathcal{I}) \nonumber \\
t &= \text{CLIP}_{text}(\mathcal{T})
\end{align}

Hệ thống tuân thủ nghiêm ngặt **3 Cấp độ Phân tách (3-Level Disentanglement)**:
\textbf{Cấp độ 1:} Chiết xuất phương thức chung và cục bộ thông qua \texttt{ModalityProjector}. Mỗi véc-tơ $(v, t)$ đi qua 2 luồng MLP để tách phần cốt lõi (general) và phần đặc thù (specific):
\begin{align}
z_{img\_g}, z_{img\_s} &= \text{ModalityProj}_{Img}(v) \\
z_{txt\_g}, z_{txt\_s} &= \text{ModalityProj}_{Txt}(t)
\end{align}

\textbf{Cấp độ 2:} Dung hợp véc-tơ chung (General Fusion):
\begin{equation}
z_{unified} = \text{Concat}(z_{img\_g}, z_{txt\_g})
\end{equation}
Ngoài ra, để gia tăng độ mạnh bạo và trơn tru cho bề mặt quyết định, hệ thống kết hợp kỹ thuật Data Augmentation dạng \textbf{MixUp} tại không gian hợp nhất $z_{unified}$, tạo ra các vector trung gian nâng cao khả năng chống chịu nhiễu.

\textbf{Cấp độ 3:} Phân tách Đặc trưng Nhân quả. Từ $z_{unified}$, mô hình áp dụng \texttt{CausalDisentanglerV2} để chia thông tin thành hai không gian trực giao:
\begin{align}
X_c, X_s &= \text{CausalDisentangle}(z_{unified})
\end{align}
$X_c$ chứa đặc trưng nhân quả quyết định nhãn, trong khi $X_s$ bắt giữ phân phối tương quan ngoại biên (spurious context/bias). Cấu trúc này thiết lập nền móng cho SCM (Mô hình Causal Cấu trúc).

\subsection{Tương tác qua Mạng Tích chập Đồ thị (GraphSAGE)}
Khác biệt với SOTA sử dụng GCN tiêu chuẩn, kiến trúc v2 triển khai Mạng \textit{GraphSAGE (Graph Sample and Aggregate)} để gia tăng khả năng bao hàm thông tin láng giềng. Bằng việc xây dựng ma trận kề $\tilde{A}$ (k-NN Graph) từ khoảng cách của véc-tơ $X_c$, thông điệp được tổ hợp (concatenate) thay vì chỉ tính tổng:
\begin{align}
h_N &= \tilde{A} X_c \\
X_c^{graph} &= \text{ReLU}\left( W_{sage} [X_c || h_N] \right)
\end{align}

\textbf{Cải tiến Động năng Cạnh (Heterophily-Aware Graph Rewiring):}
Nhược điểm của k-NN graph là dễ vướng phải kết nối giữa các node không cùng nhãn (dị biệt - Heterophily), đặc biệt là trong Unseen Domain. Để khắc phục, hệ thống giới thiệu module \texttt{EdgeHeterophilyScorer}. Dựa trên phân phối logits bước đầu, module trừng phạt các cạnh nối có mức độ dị biệt cao:
\begin{equation}
\text{Sim} = \text{Sim} + \tau \cdot \log(H\_score)
\end{equation}
Trọng số này tự động bóp nghẹt các liên kết rác, đảm bảo mạng GNN giao tiếp trên một cấu trúc đồ thị đồng nhất và sạch sẽ tuyệt đối.

\subsection{Cơ chế Điều chỉnh Cửa sau (Backdoor Adjustment)}
Để xóa bỏ "Đường dẫn cửa sau" do các yếu tố nhiễu (Confounders) gây ra, Causal GNN v2 duy trì một \textbf{Spurious Memory Bank} dung lượng 2000 phần tử (chứa $X_s$ từ các batch trước). Trong quá trình Đánh giá (Inference), định lý Backdoor Adjustment được mô phỏng bằng phép lấy mẫu ngẫu nhiên (Monte Carlo Sampling) với $N=50$:
\begin{equation}
\text{Logits}_{ba} = \frac{1}{N} \sum_{i=1}^{N} \text{Classifier}(X_c^{graph} + x_s^{(i)})
\end{equation}
Việc ghép ghép Nhân quả ($X_c^{graph}$) của mẫu hiện tại với các bối cảnh $x_s$ bất kỳ hoàn toàn trung hòa ảnh hưởng của bối cảnh, trả về phân phối xác suất thuần túy.

\section{Chiến lược Huấn luyện (Training Pipeline)}

Hệ thống được thiết kế theo tư tưởng \textbf{Curriculum Learning} nhằm định hướng sự hội tụ của không gian Causal và Graph một cách an toàn.

\subsection{Hệ thống Hàm Mất Mát (Loss Functions)}
Sự phân tách $X_c$ và $X_s$ được giám sát thông qua hàm mất mát tổng hợp của 4 luồng đạo hàm chính:
\begin{equation}
\mathcal{L}_{Total} = \alpha_1 \mathcal{L}_{Task} + \alpha_2 \mathcal{L}_{SupCon} + \alpha_3 \mathcal{L}_{Orth} + \alpha_4 \mathcal{L}_{Graph}
\end{equation}
Trong đó:
\begin{itemize}
    \item $\mathcal{L}_{Task}$ đào tạo nhánh Nhân quả bằng Cross-Entropy (sau Graph Layers).
    \item $\mathcal{L}_{SupCon}$ (Supervised Contrastive Loss) học ranh giới phương thức chặt chẽ.
    \item $\mathcal{L}_{Orth}$ (Orthogonal Loss) giữ tính trực giao cho không gian $X_c$ và $X_s$.
\end{itemize}
Đặc biệt, Mạng đối nghịch Gradient Reversal Layer (GRL) được cài cắm ở nhánh Classifier của $X_c$ để đánh lừa Domain Discriminator, chính thức ép buộc điều kiện cực kỳ khắt khe: $(X_c \perp \text{Domain})$.

\subsection{Hai Giai Đoạn Đào Tạo (Two-Phase Curriculum)}
\begin{itemize}
    \item \textbf{Phase 1: Warm-up (Epoch 1 $\rightarrow$ 15):} 
    Module GCN bị ngắt kích hoạt (Adj = None). Các bộ chiếu đa phương thức (Projections) và MLP Extractor được học cục bộ để thiết lập không gian đặc trưng cơ bản. Tại pha này, điểm dữ liệu là độc lập (I.I.D).
    
    \item \textbf{Phase 2: Graph-enhanced Learning (Epoch 16 $\rightarrow$ 40):} 
    GCN được bật qua ma trận kề Batch-wise. Đồ thị bắt đầu lan truyền thông tin (Message Passing). Đồng thời, hệ thống thu thập $X_s$ sinh ra vào \texttt{MemoryBank}. Khâu đánh giá (Evaluate) bắt đầu kích hoạt triệt để Backdoor Adjustment lấy mẫu Monte Carlo thay vì Forward chuẩn.
\end{itemize}

\section{Kết quả Thực nghiệm (Experimental Results)}

Quá trình dự đoán được thử nghiệm trên bộ CrisisMMD với 3 tác vụ chính, chạy bằng nhiều Seed (42, 2024, 2025). Kết quả Multi-Seed trung bình như sau:

\begin{table}[H]
    \centering
    \caption{Kết quả Phân loại cuối cùng (Causal GNN v2)}
    \begin{tabular}{llc}
        \toprule
        \textbf{Phân loại Tác vụ} & \textbf{Metric} & \textbf{Mean $\pm$ Std} \\
        \midrule
        \multirow{2}{*}{\textbf{Task 1: Tính Hữu ích (Informative)}} & Weighted F1 & $0.8223 \pm 0.0021$ \\
                                                                    & Balanced Acc & $0.8193 \pm 0.0022$ \\
        \midrule
        \multirow{2}{*}{\textbf{Task 2: Phân loại Nhân đạo (Humanitarian)}}& Weighted F1 & $0.6718 \pm 0.0025$ \\
                                                                    & Balanced Acc & $0.6652 \pm 0.0020$ \\
        \midrule
        \multirow{2}{*}{\textbf{Task 3: Thiệt hại Vật chất (Damage)}}    & Weighted F1 & $0.6931 \pm 0.0027$ \\
                                                                    & Balanced Acc & $0.6881 \pm 0.0000$ \\
        \bottomrule
    \end{tabular}
\end{table}

Phân tích định tính hình thái dữ liệu qua các đồ thị \textbf{t-SNE} cho thấy không gian đặc trưng của $X_c^{graph}$ trong Phase 2 đã phân cụm rõ ràng và chia tách tốt hơn hẳn so với Phase 1 không có đồ thị, khẳng định sức mạnh của Graph Reasoning khi giải quyết tính mơ hồ trong dữ liệu mạng xã hội.

\section{Kết luận (Conclusion)}
Kiến trúc Causal GNN v2 giải quyết triệt để rào cản của hiện tượng vướng mắc bối cảnh (contextual bias entanglement) trên dữ liệu thảm họa. Việc sử dụng Mạng Tích chập Đồ thị (GCN) cùng thuật toán phân tách nhân quả với Điều chỉnh Cửa sau (Backdoor Adjustment) đã chứng minh được tính hiệu quả và bền vững qua kết quả thử nghiệm định lượng lẫn định tính. Mô hình sẵn sàng phục vụ làm xương sống (backbone) phân loại cho các ứng dụng giảm nhẹ rủi ro thiên tai.

\end{document}
